\documentclass[../main.tex]{subfiles}

\begin{document}
  \section{Global Implications}
  Recall that we use $g$ to denote an arbitrary factor of $f$ and $L$ to denote the field $\Q\left[ x \right] / \left( g \right)$.
  Recall that we use $N$ to denote the product of all the inertia groups of $L$ over $\Q$ above primes not dividing $ab$.
  Corollary \ref{cor:inertia_group_acts_as_transposition} allows us to prove the following corollary:
  \begin{corollary}
    \label{cor:transitive_action}
    If $N$ acts transitively on $\Hom_{\Q}\left( L, \overline{\Q} \right)$, then the Galois group of $g$ over $\Q$ is $S_{\deg g}$.
  \end{corollary}
  \begin{proof}
    It follows from Corollary \ref{cor:inertia_group_acts_as_transposition} and the fact that any transitive permutation group generated by transpositions is symmetric (see for example \cite[Lemma 5]{Osada1987TheGG}).
  \end{proof}

  That corollary in turns allows us to prove the abstract sufficient condition for the Galois group of $g$ to be symmetric:
  \begin{thm}
    \label{thm:strategy}
    If $\Q$ is the only subextension of $L$ that is unramified away from $ab$, then the Galois group of $g$ over $\Q$ is $S_{\deg g}$.
  \end{thm}
  \begin{proof}
    By Corollary \ref{cor:transitive_action}, it suffices to show that $N$ acts transitively on $\Hom_{\Q}\left( L, \overline{\Q} \right)$.
    We will show it using the famous fact that the functor $\Hom\left( -, \overline{\Q} \right)$ is a contravariant equivalence between the category of \'{e}tale $\Q$-algebras and the category of finite $G_\Q$-sets.
    Since $N$ is normal in $G_\Q$, $G_\Q$ acts on $\Hom_{\Q}\left( L, \overline{\Q} \right) / N$  through $G_\Q/N$.
    In particular, the quotient map $\Hom_{\Q}\left( L, \overline{\Q} \right) \to \Hom_{\Q}\left( L, \overline{\Q} \right) / N$ is an epimorphism in the category of $G_\Q$-sets.
    Therefore, there is a corresponding embedding of fields $K \hookrightarrow L$.
    
    However, since $N$ acts trivially on $\Hom_{\Q}\left( L, \overline{\Q} \right) / N$, it follows that $K$ is fixed by $N$ and is hence ramified only above primes that divide $ab$.
    Therefore, $K = \Q$ and hence 
    \[
      \Hom_\Q\left( L, \overline{\Q} \right) / N \simeq 
      \Hom_\Q\left( \Q, \overline{\Q} \right) 
    \] 
    and is hence a singleton implying that $N$ indeed acts transitively on $\Hom_{\Q}\left( L, \overline{\Q} \right)$.
  \end{proof}

  Recall that we use $\rd\left( K \right)$ to denote the root discriminant of a number field $K$.
  We can combine this result with the bound on the discriminant to get an effective sufficient condition for the Galois group of $L$ to be the full symmetric group:
  \begin{thm}
    \label{thm:criterion}
    Assume that there is a constant $C_p > v_p\left( ab \right)$ for each prime $p|ab$, such that $L$ has no subextension $K$ satisfying both of the following conditions: 
    \begin{enumerate}
      \item $v_p\left( \rd\left( K \right) \right) < C_p$ for each $p|ab$,
      \item $K$ is ramified only above primes dividing $ab$.
    \end{enumerate} 
    Assume that
    \[
      \deg g > \max_{p| ab}\left( \frac{n v_p\left( ab \right) + \exrootscountbound}{C_p} \right)
    \] 
    Then the Galois group of $g$ over $\Q$ is $S_{\deg L}$.
  \end{thm}
  \begin{proof}
    We will show that the conditions of Corollary \ref{thm:strategy} hold.
    Let $K$ be a number field ramified only above primes dividing $ab$.
    Assume towards a contradiction that $K \subseteq L$.
    The inequality in the statement of the theorem is equivalent to the statement that 
    \[
      C_p > \frac{n v_p\left( ab \right) + \exrootscountbound}{\deg L}
    \]
    for each $p|ab$.
    It follows, by the contrapositive of the first assumption in the theorem statement, that there is a prime $p|ab$ such that 
    \[
      C_p \leq v_p\left( \rd\left( K \right) \right) 
    \] 
    Therefore, if we denote by $d\left( L/K \right)$ the relative discriminant of $L$ over $K$, and by $N$ the norm from $K$ to $\Q$, then
    \begin{align*}
      & v_p\left( d\left( L \right) \right) = 
      v_p\left( \d\left( K \right)^{\left[ L:K \right]} N\left( \d\left( L/K \right) \right) \right) \geq \deg L \cdot v_p\left( \rd\left( K \right) \right) \geq \\
      & \deg L \cdot C_p > \frac{n v_p\left( ab \right) + \exrootscountbound}{\deg L} = \\
      & \deg L \cdot v_p\left( ab \right) + \exrootscountbound
    \end{align*}
    which is a contradiction to proposition \ref{cor:discriminant_bound}.
  \end{proof}
\end{document}